\documentclass[10pt,a4paper]{amsart}
\usepackage[margin=19mm]{geometry}
\usepackage{amsmath,amssymb,mathtools}
\usepackage[scr=rsfs,cal=euler]{mathalfa}
\usepackage{xfrac}
\usepackage[unicode,hidelinks,bookmarks=false]{hyperref}
\newcommand{\ie}{i.e.\ }
\newcommand{\cf}{cf.\ }
\newcommand{\resp}{respectively\ }
\newcommand{\Subsection}{\subsection}
\providecommand{\nobreakdash}{\nobreak\hbox{-}}
\setlength{\emergencystretch}{2em}
\multlinegap0pt
\usepackage{bm}
\usepackage{bbm}
\usepackage{dsfont}
\usepackage{xspace}
\usepackage{cancel}
\usepackage{mathtools}
\usepackage{amsthm}
\makeatletter
\@ifundefined{lemma}{\newtheorem{lemma}{Lemma}}{}
\makeatother
\makeatletter
\newcommand{\R}{\mathbb{R}}
\expandafter\ifx\csname assumption\endcsname\relax
\newenvironment{assumption}[1][Assumption]{\begin{trivlist}
\item[\hskip \labelsep {\bfseries #1}]}{\end{trivlist}}
\fi
\newcommand{\rmd}{{\rm d}}
\newcommand{\ve}{\varepsilon}
\makeatother
\title[On the fragility of laminar flow]{On the fragility of laminar flow}
\author{Theodore D. Drivas, Daniel Ginsberg, Marc Nualart}
\date{Source: arXiv:2505.17817v1 (CC BY 4.0)}
\begin{document}
\maketitle

\noindent\emph{Source and licence.} On the fragility of laminar flow. Version arXiv:2505.17817v1 (CC BY 4.0), distributed under the Creative Commons Attribution 4.0 International licence, as recorded in the arXiv licence metadata for this exact version and shown on the abstract page https://arxiv.org/abs/2505.17817v1. Full rights notice: licenses/arxiv-2505.17817-CC-BY-4.0.txt.\par\smallskip

\noindent\textit{Editorial scope.} Counted unit: the Lemma whose source label is \texttt{lemma:ApproxGammave} in the article (source0.tex lines 312--318, source label \texttt{lemma:ApproxGammave}), delivered together with the article's own complete original proof of it (source0.tex lines 320--354, one \texttt{proof} environment of 35 lines). No mathematics is written, repaired or re-derived: statement and proof are the article's own text line for line. Required context retained uncounted: eq:defPsive (lines 263--265). Every definition, display and label of those lines is retained unchanged. Omitted from this package and not redistributed: 1--262, 266--311, 319--319, 355--1027. Nothing inside a retained excerpt is omitted or altered: every retained body is delivered exactly as the article's lines are written, so no omission receipt of any kind accompanies this package. Citations: the retained excerpts carry no citation command at all, so no bibliography entry of the article is transcribed and no reference is redistributed. Reference adaptation: none was needed and none was made; every reference inside the delivered unit resolves inside it, so no unresolved reference or citation is produced. Printed slips kept literal: no printed word, symbol, hypothesis or number of the counted statement or of its proof was altered. Layout and class: the article's own class is \texttt{amsproc} with packages including geometry, babel, inputenc, ulem, amsmath, amssymb, amsthm, comment, mathtools, hyperref, mathrsfs, xcolor, enumitem, tikz; the wrapper is the lane's pinned \texttt{amsart} editorial wrapper at 10pt on A4, which restates the theorem-like environments the retained text needs and adds the article's own macro definitions that the retained text needs (\texttt{\textbackslash R}, \texttt{\textbackslash rmd}, \texttt{\textbackslash ve}), with extra wrapper packages (bm, bbm, dsfont, xspace, cancel, mathtools). The delivered numbering is the unit's own. Assets: no retained span contains a figure environment, a graphics-inclusion command, a PDF-page inclusion, a table or a TikZ picture; no image file is copied into this package, the package declares no asset dependency and no image file is redistributed with it. Licence: version arXiv:2505.17817v1 of this article is distributed under the Creative Commons Attribution 4.0 International licence, as recorded in the arXiv licence metadata for this exact version and shown on the abstract page https://arxiv.org/abs/2505.17817v1. A full rights notice is delivered at licenses/arxiv-2505.17817-CC-BY-4.0.txt. Characters: every retained excerpt is pure ASCII, so the delivered engine, XeLaTeX with its default Latin Modern text font, reports no missing character.
\begin{equation}\label{eq:defPsive}
\psi_\ve(x,y) = {\psi_0}(x,y) + \ve {\varphi}(x,y) + r_\ve(x,y).
\end{equation}

\begin{lemma}\label{lemma:ApproxGammave}
For $\ve>0$ small enough, any singular streamline $\Gamma_\ve$ of $\psi_\ve$ is given by a graph in $x$, that is, there exists some $y_\ve\in C^1(\mathbb{T})$ such that
\begin{align*}
\Gamma_\ve = \lbrace (x,y_\ve(x)): \, x\in \mathbb{T} \rbrace.
\end{align*}
Moreover, there holds $\Vert y_\ve - y_0 \Vert_{C^1(\mathbb{T})}\lesssim \ve$.
\end{lemma}

\begin{proof}
We begin by recalling that $F(c_0)=\Delta\psi_0|_{\Gamma_0}= \partial_y^2(x,y_0(x))\left( 1 + (y_0'(x))^2\right)\neq 0$ and assume without loss of generality that $F_0:=F(c_0)<0$, so that $\partial_y^2\psi_0(x,y_0(x))<0$ for all $x\in \mathbb{T}$. By uniform continuity, we have that $\partial_y\psi_0(x,y)\leq \frac{F_0}{2}$ for all $x\in \mathbb{T}$ and all $|y-y_0(x)|<\delta_0$, for some $\delta_0>0$ sufficiently small. In particular, since
\begin{align*}
\partial_y\psi_0(x,y_0(x) + \delta) = \delta\int_0^1 \partial_y^2\psi_0(x,y_0(x) + s\delta)\rmd s ,
\end{align*} 
for all $\delta\in \R$, we conclude that $\partial_y\psi_0(x,y_0(x)+\delta) \leq \delta \frac{F_0}{2} <0$ and $\partial_y^2\psi_0(x,y_0(x) - \delta) > -\delta\frac{F_0}{2}$, for all $0 < \delta < \delta_0$. Hence, for any fixed $\delta\in (0,\delta_0)$ we observe that
\begin{align*}
\partial_y\psi_{\ve}(x,y_0(x) +\delta) &= \partial_y\psi_0(x,y_0(x) +\delta) +\left( \ve\partial_y\varphi + \partial_y r_{\ve} \right)(x,y_0(x) +\delta) \\&
\leq \delta\frac{ F_0}{2} - \ve \Vert \varphi\Vert_{C^1(D_\ve)} - \Vert r_{\ve} \Vert_{C^1(D_\ve)} <0
\end{align*}
for $\ve>0$ sufficiently small, 
where $r_\ve$ is a remainder defined in \eqref{eq:defPsive}. Likewise, we have $\partial_y \psi_{\ve}(x,y_0(x) -\delta) >0$. In particular, this shows that for all $\ve>0$ small enough, for all $x\in\mathbb{T}$, there exists $y_{\ve}(x)\in (y_0(x)-\delta, y_0(x) + \delta)$ such that $\partial_y\psi_{\ve}(x,y_{\ve}(x))=0$. Since $\partial_y^2\psi_{\ve}(x,y) =\partial_y^2\psi_0(x,y_0(x)) +O(\ve) < 0$ for all $x\in \mathbb{T}$ and all $|y-y_0(x)|<\delta_0$, for $\ve$ sufficiently small, by the implicit function theorem we conclude that the mapping $x\mapsto (x,y_{\ve}(x))$ is a unique, $2\pi$-periodic and $C^1$ curve, with $\Vert y_\ve(x) - y_0(x) \Vert_{C(\mathbb{T})}\leq \delta$.

In fact, once we know $\Vert y_\ve(x) - y_0(x) \Vert_{C(\mathbb{T})}\leq \delta$ for all $\delta\in (0,\delta_0)$ by taking $\ve>0$ small enough, we can improve it to $\Vert y_\ve(x) - y_0(x) \Vert_{C^1(\mathbb{T})}\lesssim \ve$. Indeed, for all $x\in \mathbb{T}$ we have that
\begin{align}\label{eq:partialypsixyve}
0 &= \partial_y\psi_\ve(x,y_\ve(x)) = \partial_y \psi_0(x,y_\ve(x)) + \ve\partial_y\varphi(x,y_\ve(x)) + \partial_y r_\ve(x,y_\ve(x)) 
\end{align}
where further
\begin{align}\label{eq:partialypsi0xyve}
\partial_y \psi_0(x,y_\ve(x)) = (y_\ve(x) - y_0(x) ) \int_0^1 (\partial_y^2\psi_0)(x,y_0(x) + t(y_\ve(x) - y_0(x))) \rmd t.
\end{align}
Since $\partial_y^2\psi_0(x,y_0(x))\neq 0$, by uniform continuity we have that $\partial_y^2\psi_0(x,y)\neq 0$ for all $|y-y_0(x)|\leq \delta_0$, for some $\delta_0>0$. For $\ve$ small enough such that $\Vert y_\ve(x) - y_0(x) \Vert_{C(\mathbb{T})}\leq \delta_0$, we then have
\begin{align*}
\Vert y_\ve(x) - y_0(x) \Vert_{C(\mathbb{T})}\lesssim \ve \Vert \varphi \Vert_{C^1(D_{\ve})} + \Vert r_\ve \Vert_{C^1(D_{\ve})} \lesssim \ve.
\end{align*}
Similarly, applying $\frac{\rmd}{\rmd x}$ to \eqref{eq:partialypsixyve}  and \eqref{eq:partialypsi0xyve} we have that
\begin{align*}
0 &= (y_\ve(x) - y_0(x) )' \int_0^1 (\partial_y^2\psi_0)(x,y_0(x) + t(y_\ve(x) - y_0(x))) \rmd t \\  
&\quad + (y_\ve(x) - y_0(x) )'(y_\ve(x) - y_0(x) ) \int_0^1 (\partial_y^2\psi_0)(x,y_0(x) + t(y_\ve(x) - y_0(x))) \rmd t \\
&\quad + (y_\ve(x) - y_0(x) ) \int_0^1 (\partial_{xyy}^3\psi_0 + (y_0'(x))\partial_y^3\psi_0)(x,y_0(x) + t(y_\ve(x) - y_0(x))) \rmd t \\
&\quad + (y_\ve(x)-y_0(x))' \left( \ve \partial_y\varphi(x,y_\ve(x)) + \partial_y r_\ve(x,y_\ve(x)) \right) \\
&\quad + (\partial_x + y_0'(x)\partial_y)\left(\ve \partial_y\varphi(x,y_\ve(x)) + \partial_y r_\ve(x,y_\ve(x)) \right).
\end{align*}
Since $\Vert y_\ve - y_0 \Vert_{C(\mathbb{T})}\lesssim \ve$ and $\ve \Vert \varphi \Vert_{C^2(D_{\ve})} + \Vert r_\ve \Vert_{C^2(D_{\ve})}\lesssim \ve$, taking $\ve$ small enough we conclude that $\Vert y_\ve - y_0 \Vert_{C^1(\mathbb{T})}\lesssim \ve$. Finally, let $\Gamma_{\ve} := \lbrace (x, y_{\ve}(x)) : x\in \mathbb{T} \rbrace$ and note that $\partial_y\psi_{\ve}|_{\Gamma_{\ve}}=0$ and $\psi_\ve$ is periodic on $\Gamma_{\ve}$. Hence, there exists $x_0\in \mathbb{T}$ for which $\nabla\psi_{\ve}(x_0,y_{\ve}(x_0))=0$. If $\psi_{\ve}$ is laminar, then $\nabla\psi_{\ve}=0$ on $\Gamma_{\ve}$.
\end{proof}

\end{document}
